\documentclass[a4paper]{article}

\usepackage[polish]{babel}
\usepackage[T1]{fontenc}
\usepackage[utf8]{inputenc}
\usepackage{graphicx}
\usepackage{fancyhdr}
\pagestyle{fancy}

\lhead[]{}
\chead[Fundacja CoderDojo Polska]{Fundacja CoderDojo Polska}
\rhead[]{}
\cfoot[Treść udostępniana na licencji Creative Commons – Uznanie Autorstwa – Na Tych Samych Warunkach]{Treść udostępniana na licencji Creative Commons – Uznanie Autorstwa – Na Tych Samych Warunkach}

\title{Jak założyć nowy klub CoderDojo?}
\author{Fundacja CoderDojo Polska}
\date{\today}


\begin{document}
\maketitle
\includegraphics[width=\textwidth]{logo_CD.pdf}

\section{Kroki}
\begin{enumerate}
\item Znajdź najbliższe dojo na naszej mapie (link). 
\item W miarę możliwości przyjdź do nas na zajęcia i porozmawiaj z mentorami. Jeśli masz do nas za daleko skontaktuj się z nimi np. na Skype. Chętnie opowiedzą Ci co i jak.  
\item Znajdź drugiego mentora. Może to być osoba znająca się na programowaniu, ale nie musi. Fajni są też animatorzy, pedagodzy i każdy, kto chce się zaangażować w działania CoderDojo. 
\item Znajdź salę w której będą odbywały się spotkania* Jeśli masz z tym problem – napisz do nas. Chętnie Ci pomożemy!
\item Poproś nas o materiały. Mamy scenariusze, samouczki i dużą kolektywną wiedzę na tematy związane z robotyką i programowaniem. W CoderDojo uczymy się od siebie nawzajem. 
\end{enumerate}


\section{Co musi mieć założyciel CoderDojo}
\begin{enumerate}
\item Determinację. Chcielibyśmy, żeby Ci się udało, a gwarancją sukcesu jest chęć. 
\item Czas, aby zaangażować się w działanie dojo. Wszyscy mentorzy działają na zasadzie wolontariatu. Musisz więc mieć pewność, że masz czas i możliwość zająć się tworzeniem nowego dojo.   
\item Ludzi, którzy pomogą mu prowadzić zajęcia. Rodzice uczestniczących dzieci też są przydatni i ważni. 
\end{enumerate}

\section{Czego nie musi mieć założyciel CoderDojo}
\begin{itemize}
\item Pozwolenia na działanie, uprawnień i dyplomów. Każdy może założyć własne dojo. 
\item Sprzętu dla dzieci. W CoderDojo hołdujemy zasadzie BYOD (\textit{Bring Your Own Device}). Jeśli gdzieś są  dzieci, które nie mają własnych komputerów mamy możliwości, by wypożyczyć im sprzęt. 
\item Grupy chętnych dzieci. Nie uzależniaj decyzji o założeniu dojo od tego, czy masz chętnych. Z Naszego doświadczenia wynika, że jak zaczniesz działać dzieci się znajdą. 
\end{itemize}

\section{Idealna sala na CoderDojo}
\begin{itemize}
\item Powinna znajdować się w MOKu, siedzibie lokalnej firmy, w Techno Parku lub bibliotece. Unikajmy szkół – dzieciom kojarzą się z przymusem, a nie mają zazwyczaj odpowiedniej infrastruktury na dojo.  
\item Musi mieć podpięcie do internetu i dobre WiFi
\item Ważna jest lokalizacja. Pamiętaj, że co sobotę ty i uczestnicy będziecie musieli do niej dojeżdżać. Sprawdź komunikację miejską i możliwość parkowania. 
\item Powinna być bezpieczna. Unikajmy miejsc, gdzie kręci się dużo postronnych osób, gdzie każdemu łatwo się „wmieszać” w zajęcia. 
\end{itemize}







\end{document}